\chapter{Twin Primes and Prime Gaps}\label{ch:twins}

\begin{center}\itshape
Every even number is the difference of two consecutive primes in infinitely many ways.
\upshape(Alphonse de Polignac, 1849, in our words)
\end{center}

On 17 April 2013 a paper arrived at the \emph{Annals of Mathematics} from Yitang Zhang, a lecturer at the University of New Hampshire who was then
almost unknown. It proved that there are infinitely many pairs of primes that differ by at most $70\,000\,000$. Seventy million is a long way from $2$,
but it is finite, and nobody had ever proved that any finite gap occurs infinitely often. Within a year an online collaboration, the Polymath8 project,
brought the bound down to $4680$, and James Maynard, then a postdoctoral researcher, found a simpler and stronger method that gave $600$; combined with
Polymath, the bound became $246$. The twin prime conjecture asks for $2$.

This chapter is about the gaps between consecutive primes $p_n<p_{n+1}$, written $d_n=p_{n+1}-p_n$. On average they grow like $\log p_n$, by the prime
number theorem. The questions are about how far they stray from the average: how small can they be, infinitely often, and how large? We meet the
Hardy--Littlewood conjecture that predicts how many twins, cousins and triples there are, with constants that match a billion numbers to four
decimals. We prove two classical theorems at the two ends: Erd\H{o}s's theorem of 1940 that gaps noticeably smaller than average occur infinitely often,
and Westzynthius's theorem of 1931 that gaps can be arbitrarily many times larger than average. And we explain why the sieve, the tool that proves
both, cannot reach the twin primes on its own.

\begin{think}
Near $10^9$, how large is a typical gap between consecutive primes? Can $p$, $p+2$ and $p+4$ all be prime for some $p>3$? And $p$, $p+2$, $p+6$?
\end{think}

\section{Patterns of primes}

Primes come in patterns. Twins $p,p+2$ (like $101,103$), cousins $p,p+4$ ($103,107$), sexy primes $p,p+6$, and triples $p,p+2,p+6$ ($101,103,107$).
Not every pattern is possible: among $p$, $p+2$, $p+4$ one is divisible by $3$, so only $3,5,7$ has all three prime. A finite set
$H=\{h_1,\dots,h_k\}$ of integers is \emph{admissible} if for every prime $p$ the $h_i$ miss at least one residue class modulo $p$; equivalently,
if $\nu_H(p)<p$ for every $p$, where $\nu_H(p)$ is the number of classes the $h_i$ occupy. Only primes $p\le k$ need checking. If $H$ is not
admissible, some prime divides one of $n+h_1,\dots,n+h_k$ for every $n$, and the pattern occurs only finitely often.

\begin{remark}[Hardy--Littlewood, 1923]
The \emph{prime $k$-tuples conjecture} says that every admissible $H$ occurs infinitely often, and more precisely that the number of $n\le x$ with
$n+h_1,\dots,n+h_k$ all prime is asymptotic to
\[
\mathfrak S(H)\int_2^x\frac{\mathrm dt}{(\log t)^k},\qquad\mathfrak S(H)=\prod_p\Bigl(1-\frac{\nu_H(p)}p\Bigr)\Bigl(1-\frac1p\Bigr)^{-k}.
\]
The product converges, and is positive exactly when $H$ is admissible (Tome~5, Proposition~10.5). Tome~5, Chapter~10 derives the formula from
Cram\'er's model corrected prime by prime. For twins and cousins $\mathfrak S=2C_2=1.3203\ldots$, with the twin prime constant
$C_2=\prod_{p>2}\bigl(1-\frac1{(p-1)^2}\bigr)=0.6601618\ldots$; for the triple $\{0,2,6\}$, $\mathfrak S=\frac92\prod_{p\ge5}\bigl(1-\frac{3p-1}{(p-1)^3}\bigr)=2.858248\ldots$
\end{remark}

How good is the prediction? We sieved the integers up to $10^9$ and counted (Figure~\ref{fig:tuples}). There are $3\,424\,506$ twin pairs with
$p\le10^9$, and the prediction is $3\,425\,308$, a ratio of $0.99977$. Cousins: $3\,424\,680$ against the same prediction. Triples $p,p+2,p+6$: $379\,508$ against
$379\,803$, a ratio of $0.99922$. Predictions this precise, with constants computed from a product over all primes, are why number theorists believe
the conjecture although nobody can prove that even one pattern occurs infinitely often.

\begin{figure}[ht]
\centering
\fig{ch02-tuples}
\caption{The number of $p\le x$ such that $p,p+2$ (blue), $p,p+4$ (green) and $p,p+2,p+6$ (red) are all prime, divided by the Hardy--Littlewood
prediction, for $10^2\le x\le10^9$. All three ratios tend to $1$; at $10^9$ they are $0.99977$, $0.99982$ and $0.99922$.}
\label{fig:tuples}
\end{figure}

Admissible sets are easy to find, and they can be as dense as the primes themselves.

\begin{proposition}\label{prop:firstprimes}
For every $k\ge1$, the first $k$ primes larger than $k$ form an admissible set.
\end{proposition}
\begin{proof}
Let $H$ be this set. For a prime $p\le k$, no element of $H$ is divisible by $p$ (they are primes larger than $k$), so the class $0$ is missed. For
$p>k$, the $k$ elements occupy at most $k<p$ classes.
\end{proof}

By the prime number theorem, the diameter of this set is about $k\log k$. This is how the bound $246$ arises: Maynard's method shows that every
admissible set of $k$ elements, for a suitable $k$ ($k=105$ in Maynard's paper, $k=50$ after Polymath), contains at least two primes in infinitely many
translates, and the narrowest admissible set of $50$ elements has diameter $246$.

\section{Erd\H{o}s: small gaps occur}

By the prime number theorem, the $\pi(x)-\pi(x/2)\sim\frac x{2\log x}$ primes in $]\frac x2,x]$ share a total gap of at most $\frac x2$, so their average
gap is about $\log x$. Some gaps must be at most average, so $\liminf\frac{d_n}{\log p_n}\le1$ (Exercise~\ref{ex:averagegap}). To do better, we must
show that the gaps cannot all bunch up just below the average, and for that we need an upper bound for the number of prime pairs with a given
difference. This is Brun's theorem for twins (Tome~5, Theorem~8.5), redone for every even difference $h$, with a bound uniform in $h$.

\begin{proposition}\label{prop:pairs}
There is an absolute constant $C$ such that for every even $h\ge2$ and every $x\ge3$,
\[
\pi_h(x)=\#\{p\le x:\ p+h\text{ prime}\}\le C\Bigl(\frac h{\varphi(h)}\Bigr)^2\frac x{(\log x)^2}.
\]
\end{proposition}
\begin{proof}
We use Selberg's sieve (Tome~5, Theorem~8.2) with its notation, writing $\eta$ for the function called $h$ there, to avoid a clash with our $h$.

\emph{The data.} Let $\mathcal A=\{n(n+h):1\le n\le x\}$, $\mathcal P$ all primes, and $y=z=\frac{\sqrt x}{(\log x)^3}$. If $p>z$ and $p+h$ are both prime, then
$n=p$ is counted by $S(\mathcal A,\mathcal P,z)$, so $\pi_h(x)\le z+S(\mathcal A,\mathcal P,z)$. For squarefree $d$, the number $\rho(d)$ of classes $n\bmod d$ with
$d\mid n(n+h)$ is multiplicative, with $\rho(p)=1$ if $p\mid h$ (the classes $0$ and $-h$ coincide) and $\rho(p)=2$ otherwise; in particular $\rho(2)=1$. So
$|\mathcal A_d|=\frac{\rho(d)}dx+r_d$ with $|r_d|\le\rho(d)$, $g(p)=\frac{\rho(p)}p$, and $\eta(p)=\frac{\rho(p)}{p-\rho(p)}$, which is $\frac2{p-2}\ge\frac2{p-1}$ for
$p\nmid h$.

\emph{The main term.} Keeping only the $d$ prime to $h$ (hence odd) in $G(y)$, and writing $2^{\omega(d)}$ as the number of ways to split $d=ab$ into
two coprime factors,
\[
G(y)\ge\sum_{\substack{d\le y\\(d,h)=1}}\frac{\mu^2(d)2^{\omega(d)}}{\varphi(d)}\ge\sum_{\substack{a\le\sqrt y\\(a,h)=1}}\frac{\mu^2(a)}{\varphi(a)}\sum_{\substack{b\le\sqrt y\\(b,ah)=1}}\frac{\mu^2(b)}{\varphi(b)}.
\]
We need two lower bounds, valid for every $q\ge1$ and $Y\ge1$:
\[
\text{(i)}\ \sum_{\substack{b\le Y\\(b,q)=1}}\frac{\mu^2(b)}{\varphi(b)}\ge\frac{\varphi(q)}q\log Y,\qquad\text{(ii)}\ \sum_{\substack{a\le Y\\(a,q)=1}}\frac{\mu^2(a)}a\ge\frac6{\pi^2}\,\frac{\varphi(q)}q\log Y.
\]
For (i), every squarefree $m\le Y$ is uniquely $m=kb$ with $k$ dividing the product $q'$ of the primes of $q$ and $(b,q)=1$; so, by Tome~5, Lemma~8.3,
$\log Y\le\sum_{m\le Y}\frac{\mu^2(m)}{\varphi(m)}\le\sum_{k\mid q'}\frac1{\varphi(k)}\cdot\sum_{b}\frac{\mu^2(b)}{\varphi(b)}$, and
$\sum_{k\mid q'}\frac1{\varphi(k)}=\prod_{p\mid q}\bigl(1+\frac1{p-1}\bigr)=\frac q{\varphi(q)}$. For (ii), the same factorisation, with $k$ now running over all
integers whose primes divide $q$, gives $\log Y\le\sum_{m\le Y}\frac1m\le\frac q{\varphi(q)}\sum_{m\le Y,(m,q)=1}\frac1m$; and every $m$ prime to $q$ is $a\ell^2$ with $a$
squarefree and prime to $q$, so $\sum_{m\le Y,(m,q)=1}\frac1m\le\frac{\pi^2}6\sum_{a\le Y,(a,q)=1}\frac{\mu^2(a)}a$.

Apply (i) to the inner sum with $q=ah$: since $(a,h)=1$, $\frac{\varphi(ah)}{ah}=\frac{\varphi(a)}a\frac{\varphi(h)}h$, so the inner sum is at least
$\frac{\varphi(a)}a\frac{\varphi(h)}h\log\sqrt y$. Then apply (ii) with $q=h$:
\[
G(y)\ge\frac{\varphi(h)}h\log\sqrt y\sum_{\substack{a\le\sqrt y\\(a,h)=1}}\frac{\mu^2(a)}a\ge\frac6{\pi^2}\Bigl(\frac{\varphi(h)}h\Bigr)^2\Bigl(\frac{\log y}2\Bigr)^2.
\]

\emph{The remainder.} As for twins, $\rho([d_1,d_2])\le2^{\omega(d_1)}2^{\omega(d_2)}\le\tau(d_1)\tau(d_2)$, so the remainder is at most
$\bigl(\sum_{d\le y}\tau(d)\bigr)^2\le y^2(1+\log y)^2\le\frac x{(\log x)^4}$ for large $x$.

\emph{Conclusion.} Since $\log y\ge\frac13\log x$ for large $x$, Selberg's theorem gives
$\pi_h(x)\le z+\frac{2\pi^2\cdot9}3\bigl(\frac h{\varphi(h)}\bigr)^2\frac x{(\log x)^2}+\frac x{(\log x)^4}$ for $x\ge x_0$, with $x_0$ independent of $h$. For
$3\le x<x_0$ use $\pi_h(x)\le x$ and enlarge $C$.
\end{proof}

The factor $\bigl(\frac h{\varphi(h)}\bigr)^2$ is an artefact of our crude lower bound for $G(y)$; the truth, according to Hardy and Littlewood, is
$\mathfrak S(\{0,h\})$, which is of the order of $\frac h{\varphi(h)}$. Either is harmless, because it is bounded on average.

\begin{lemma}\label{lem:hphi}
There are constants $C_1,C_2$ such that for all $A\ge0$ and $K\ge1$,
\[
\sum_{A<h\le A+K}\Bigl(\frac h{\varphi(h)}\Bigr)^2\le C_1K+C_2\bigl(\log(A+K+3)\bigr)^6.
\]
\end{lemma}
\begin{proof}
We have $\bigl(\frac h{\varphi(h)}\bigr)^2=\prod_{p\mid h}(1+f(p))=\sum_{d\mid h}\mu^2(d)f(d)$ with $f(p)=\bigl(\frac p{p-1}\bigr)^2-1=\frac{2p-1}{(p-1)^2}$, extended
multiplicatively. One checks that $f(p)\le\frac6p$ for every prime $p$ (this is $4p^2-11p+6\ge0$, true for $p\ge2$). Exchanging the sums, and using
that an interval of length $K$ contains at most $\frac Kd+1$ multiples of $d$,
\[
\sum_{A<h\le A+K}\Bigl(\frac h{\varphi(h)}\Bigr)^2\le\sum_{d\le A+K}\mu^2(d)f(d)\Bigl(\frac Kd+1\Bigr)\le K\prod_p\Bigl(1+\frac6{p^2}\Bigr)+\prod_{p\le A+K}\Bigl(1+\frac6p\Bigr).
\]
The first product converges. The second is at most $\exp\bigl(6\sum_{p\le X}\frac1p\bigr)$ with $X=A+K+3$, and $\sum_{p\le X}\frac1p\le\log\log X+1$ for large $X$
(Mertens, Tome~5, Theorem~1.5), which gives $e^6(\log X)^6$ (and a constant covers small $X$).
\end{proof}

\begin{theorem}[Erd\H{o}s, 1940]\label{thm:erdossmall}
There is $\delta>0$ such that $d_n\le(1-\delta)\log p_n$ for infinitely many $n$. That is, $\displaystyle\liminf_{n\to\infty}\frac{p_{n+1}-p_n}{\log p_n}<1$.
\end{theorem}
\begin{proof}
Let $C,C_1$ be the constants of Proposition~\ref{prop:pairs} and Lemma~\ref{lem:hphi}, put $B=8CC_1$ and $\delta=\frac1{8B}$, and suppose, for a
contradiction, that $d_n>(1-\delta)\log p_n$ for all large $n$. Let $x$ be large, $L=\log x$, and let $q_1<\dots<q_N$ be the primes in $]\frac x2,x]$.
Their gaps $e_i=q_{i+1}-q_i$ ($1\le i<N$) satisfy:
\begin{itemize}
\item $\sum_ie_i\le\frac x2$;
\item $e_i>(1-\delta)\log\frac x2\ge(1-2\delta)L$ for all $i$, by our assumption;
\item $N-1\ge(1-\varepsilon)\frac x{2L}$, for any fixed $\varepsilon>0$ once $x$ is large, by the prime number theorem (Tome~5, Theorem~5.5).
\end{itemize}
Let $I=[(1-2\delta)L,(1+2\delta)L]$. A gap $e_i=h\in I$ is a pair of primes $q_i\le x$, $q_i+h$, so by Proposition~\ref{prop:pairs} and
Lemma~\ref{lem:hphi}, the number of $i$ with $e_i\in I$ is at most
\[
\sum_{h\in I}\pi_h(x)\le\frac{Cx}{L^2}\Bigl(C_1(4\delta L+1)+C_2\log^6(3L)\Bigr)\le\frac{8CC_1\delta x}L=\frac{B\delta x}L
\]
for large $x$. Every other gap exceeds $(1+2\delta)L$, that is, the lower bound $(1-2\delta)L$ by at least $4\delta L$. Hence
\[
\frac x2\ge\sum_ie_i\ge(N-1)(1-2\delta)L+\Bigl(N-1-\frac{B\delta x}L\Bigr)4\delta L\ge\frac x2\Bigl((1-\varepsilon)(1+2\delta)-8B\delta^2\Bigr).
\]
With $8B\delta^2=\delta$, the bracket is $(1-\varepsilon)(1+2\delta)-\delta=1+\delta-\varepsilon(1+2\delta)$, which exceeds $1$ when $\varepsilon<\frac\delta{1+2\delta}$.
This contradiction proves the theorem.
\end{proof}

The constant $1-\delta$ was improved many times in the following decades. In 2005 Daniel Goldston, J\'anos Pintz and Cem Y\i ld\i r\i m proved that the
$\liminf$ is $0$: there are infinitely many gaps smaller than any fixed fraction of the average. Their method, a sieve with weights adapted to an
admissible tuple, was the starting point of Zhang's and Maynard's work, which is proved in a later tome.

\section{Westzynthius: large gaps occur}

At the other end, the gaps can be large. The easy construction $n!+2,n!+3,\dots,n!+n$ gives $n-1$ consecutive composite numbers, but near
$n!$ the average gap is already $\log n!\approx n\log n$, so this is far below average (Exercise~\ref{ex:factorial}). Erik Westzynthius, a student in
Helsinki, proved in 1931 that gaps can be any number of times the average.

\begin{figure}[ht]
\centering
\fig{ch02-gaps}
\caption{Maximal gaps: each blue dot is a gap $d_n$ larger than all earlier gaps, for $p_n\le10^9$ (the last is $282$, after $436\,273\,009$), plotted
against $(\log p_n)^2$; the red dot is the largest gap found below $4\cdot10^{18}$, $1476$ after $1\,425\,172\,824\,437\,699\,411$ (Oliveira e Silva, Herzog
and Pardi). The dashed line is Cram\'er's prediction $d_n\approx(\log p_n)^2$.}
\label{fig:gaps}
\end{figure}

\begin{theorem}[Westzynthius, 1931]\label{thm:westzynthius}
The ratio $\frac{p_{n+1}-p_n}{\log p_n}$ is unbounded:
\[
\limsup_{n\to\infty}\frac{p_{n+1}-p_n}{\log p_n}=\infty.
\]
\end{theorem}

The idea is to build a long run of composite numbers by the Chinese remainder theorem. If for every prime $p\le x$ we choose a residue class $a_p$,
and the classes cover the interval $[1,y]$ (every $n\le y$ is $\equiv a_p\pmod p$ for some $p\le x$), then any $N\equiv-a_p\pmod p$ for all $p\le x$ makes
$N+1,\dots,N+y$ all divisible by small primes. The product of the primes up to $x$ is about $e^x$, so $\log N$ is about $x$, and we win if $y$ is a large
multiple of $x$. Covering $[1,y]$ by one class for each prime up to $x$ is a puzzle, and the solution below wastes almost nothing: most primes
take the class $0$, the smallest primes take the class $1$, and the largest primes mop up one survivor each.

\begin{proof}
Fix $K\ge2$, and fix $z\ge2K$ (chosen below in terms of $K$ only). Let $Q=\prod_{p\le z}p$. Let $x$ be large and $y=Kx$; then $(\frac x2)^2>y$.

\emph{Step 1: classes $1$ and $0$.} Put $a_p=1$ for $p\le z$ and $a_p=0$ for $z<p\le\frac x2$. An integer $n\in[1,y]$ not yet covered has no prime factor in
$]z,\frac x2]$, so $n=st$ where $s$ has all its prime factors $\le z$ and $t$ all its prime factors $>\frac x2$. Since $(\frac x2)^2>y$, $t$ is $1$ or a prime
$q>\frac x2$. The survivors are of two kinds.
\begin{itemize}
\item $n=s$ with all prime factors $\le z$. Each exponent is at most $\log_2y$, so there are at most $(1+\log_2y)^{\pi(z)}$ of them: a power of $\log x$,
since $z$ is fixed.
\item $n=sq$ with $q>\frac x2$ prime, so $s<2K\le z$. Moreover $sq\not\equiv1\pmod p$ for every $p\le z$. For $p\mid s$ this is automatic; for $p\nmid s$ it
excludes one class of $q$ modulo $p$. So $q$ lies in a set $A_s$ of classes modulo $Q$ with
\[
\frac{|A_s|}{\varphi(Q)}=\prod_{\substack{p\le z\\p\nmid s}}\Bigl(1-\frac1{p-1}\Bigr)\le\prod_{2K<p\le z}\Bigl(1-\frac1p\Bigr)=:\varepsilon,
\]
since no prime $p>2K$ divides $s$. By the prime number theorem for progressions (Tome~5, Theorem~5.8) with the fixed modulus $Q$, the number of such
$q\le\frac ys$ is at most $(1+o(1))\varepsilon\frac{y/s}{\log(x/2)}$, and summing over $s<2K$, the number of survivors of this kind is at most
$(1+o(1))\,\varepsilon\,(1+\log2K)\frac{Kx}{\log x}$.
\end{itemize}
Since $\prod_{p\le z}(1-\frac1p)\to0$ as $z\to\infty$ (Tome~5, Theorem~1.6), we can choose $z$ so that $\varepsilon K(1+\log2K)\le\frac14$. Then, for $x$ large,
fewer than $\frac x{3\log x}$ integers of $[1,y]$ survive Step~1.

\emph{Step 2: mopping up.} There are at least $\frac x{3\log x}$ primes in $]\frac x2,x]$ for large $x$ (prime number theorem), and none has been used yet.
Give each survivor $n$ its own prime $p$ from this range and put $a_p\equiv n\pmod p$. The remaining primes get any class. Now every $n\in[1,y]$ is
$\equiv a_p\pmod p$ for some $p\le x$.

\emph{Step 3: the gap.} By the Chinese remainder theorem there is $N$ with $P<N\le2P$, where $P=\prod_{p\le x}p$, and $N\equiv-a_p\pmod p$ for every
$p\le x$. For $1\le n\le y$, some $p\le x$ divides $N+n$, and $N+n>P\ge p$, so $N+n$ is composite. Let $p_m$ be the largest prime $\le N$. Then
$p_{m+1}>N+y$, so $d_m>y=Kx$. And $\log p_m\le\log2P=\log2+\theta(x)\le\log2+x\log4$ (Tome~5, Proposition~1.2), so
\[
\frac{d_m}{\log p_m}\ge\frac{Kx}{\log2+x\log4}\ge\frac K2\qquad(x\ge2).
\]
As $x\to\infty$, $d_m>Kx\to\infty$, so this happens for infinitely many $m$. Hence the $\limsup$ is at least $\frac K2$, for every $K$.
\end{proof}

\begin{remark}[how large?]
Optimising the same construction, with smooth-number estimates and a sieve for the primes $q$, Robert Rankin (1938) obtained gaps of size
$c\log p\,\frac{\log\log p\,\log\log\log\log p}{(\log\log\log p)^2}$ infinitely often, and for 76 years only the constant $c$ improved. In 2014 two groups,
Ford, Green, Konyagin and Tao, and independently Maynard, showed that $c$ can be taken arbitrarily large, and the five authors then improved the
shape of the bound together (Journal of the AMS, 2018). Tome~14, Chapter~11 proves Rankin's bound in full, together with Maier's theorem. Cram\'er's model predicts gaps as large as $(\log p)^2$ (Tome~5, Theorem~10.4), and
Figure~\ref{fig:gaps} shows the records sitting a little below that line; the record $1476$ is $0.845(\log p)^2$. Granville argued in 1995 that the
true $\limsup$ of $\frac{d_n}{(\log p_n)^2}$ is at least $2e^{-\gamma}\approx1.12$, because the model ignores small primes.
\end{remark}

\section{Why it is hard: the parity problem}

Both theorems of this chapter were proved with a sieve, and so were Brun's theorem and the best upper bound for twins. Why does the sieve not prove
that twins exist? Selberg found the reason, now called the \emph{parity problem}. A sieve only uses the sizes $|\mathcal A_d|$ of the multiples of each
$d$. Take the integers $n\le x$ with an even number of prime factors, and those with an odd number: for every squarefree $d$, the multiples of $d$ split
almost evenly between the two sets, so the sieve sees the same data in both, up to error terms. The primes all lie in the second set. Any lower bound
the sieve proves must hold for the first set too, which contains no primes at all, so the sieve alone can never prove that a sequence contains primes.
It can prove upper bounds, and lower bounds for numbers with at most two prime factors: Chen's theorem (1973) gives infinitely many primes $p$ with
$p+2$ prime or a product of two primes. To go further, Zhang and Maynard combined the sieve with information about primes in arithmetic progressions
on average (the Bombieri--\allowbreak Vinogradov theorem), which is not sieve data, and asked for less: two primes among $k$ numbers, rather than among $2$.

\section{Reading a paper: Granville's survey}

Andrew Granville's \emph{Primes in intervals of bounded length} (Bulletin of the AMS, 2015) tells the story from Hardy and Littlewood to Maynard, and
proves Maynard's theorem in full in its middle sections.
\begin{itemize}
\item \emph{What it proves.} For every $m$ there is a bound $B_m$ such that infinitely many intervals of length $B_m$ contain $m$ primes. For $m=2$ this
gives bounded gaps.
\item \emph{The skeleton.} Fix an admissible $k$-tuple. Weigh each $n$ by the square of a sum over divisors of $(n+h_1)\cdots(n+h_k)$, as in Selberg's sieve,
but with weights that depend on how the divisors split among the factors. Compare the sum of the weights with the sum of the weights times the number
of primes among the $n+h_i$. If the ratio exceeds $m-1$, some $n$ has $m$ primes. Maynard's insight is that the weights can be chosen by optimising a
quadratic form in $k$ variables, and the ratio grows like $\log k$.
\item \emph{The idea worth stealing.} When a one-dimensional choice fails, let the weights depend on many variables and optimise: the dimension itself
pays for the loss.
\item \emph{Prerequisites.} Selberg's sieve (Tome~5, Chapter~8) and the Bombieri--\allowbreak Vinogradov theorem, which is stated in the survey and proved in a
later tome.
\end{itemize}

\section{Open doors}

\begin{itemize}
\item Count the gaps of each size $h$ up to $10^9$ and compare with the Hardy--Littlewood prediction for consecutive primes, which needs inclusion and
exclusion over the primes in between. Why are gaps of $6$ the most frequent below $10^{9}$, and which gap becomes the most frequent later?
\item Find the narrowest admissible sets of $k$ elements for $k\le20$ by computer, and compare with the first $k$ primes larger than $k$.
\item The Jacobsthal function $j(m)$ is the length of the longest run of consecutive integers sharing a factor with $m$, plus one. Compute $j$ of the
product of the first $r$ primes for $r\le10$ and compare with the bound behind Westzynthius's proof.
\item Brun's constant, the sum of $\frac1p+\frac1{p+2}$ over twin pairs, is about $1.902$. Estimate it from the twins up to $10^9$ with the
Hardy--Littlewood correction for the tail.
\end{itemize}

\section{Exercises}

\begin{exercise}[$\star$]\label{ex:admissible2}
Which of $\{0,2,4\}$, $\{0,2,6\}$, $\{0,4,6\}$, $\{0,2,6,8\}$, $\{0,2,8\}$, $\{0,6,12,18,24\}$ are admissible?
\end{exercise}

\begin{exercise}[$\star$]\label{ex:twin12}
Let $p,p+2$ be twin primes with $p>3$. Show that $12\mid2p+2$ and that $p(p+2)+1$ is a perfect square.
\end{exercise}

\begin{exercise}[$\star$]\label{ex:averagegap}
Using the prime number theorem, show that $\liminf\frac{d_n}{\log p_n}\le1\le\limsup\frac{d_n}{\log p_n}$.
\end{exercise}

\begin{exercise}[$\star$]\label{ex:factorial}
The numbers $n!+2,\dots,n!+n$ are composite. Show that the gap they produce is only about $\frac{\log p}{\log\log p}$, where $p$ is the prime before them.
\end{exercise}

\begin{exercise}[$\star$]\label{ex:diameter}
Show that the admissible set of Proposition~\ref{prop:firstprimes} has diameter $(1+o(1))k\log k$.
\end{exercise}

\begin{exercise}[$\star\star$]\label{ex:hphiavg}
Show that $\sum_{n\le x}\frac n{\varphi(n)}=\frac{\zeta(2)\zeta(3)}{\zeta(6)}x+O(\log x)$.
\end{exercise}

\begin{exercise}[$\star\star$]\label{ex:triplesieve}
Adapt the proof of Proposition~\ref{prop:pairs} to show that the number of $p\le x$ with $p$, $p+2$, $p+6$ all prime is $O\bigl(\frac x{(\log x)^3}\bigr)$.
\end{exercise}

\begin{exercise}[$\star\star$]\label{ex:jacobsthal}
Find the longest run of consecutive integers each of which shares a factor with $30$, and with $210$. Deduce, by the argument of
Theorem~\ref{thm:westzynthius}, a run of $9$ consecutive composite numbers below $420$, and find the first run of $9$ consecutive composites.
\end{exercise}

\begin{exercise}[$\star$]\label{ex:pigeon}
Assume the theorem of Maynard and Polymath: $d_n\le246$ for infinitely many $n$. Show that some even $h\le246$ is a gap $d_n$ for infinitely many $n$.
Can you say which one?
\end{exercise}

\begin{exercise}[$\star\star$]\label{ex:bruntail}
Using Brun's bound (Tome~5, Theorem~8.5), show that the sum of $\frac1p$ over the twin primes $p>x$ is $O\bigl(\frac1{\log x}\bigr)$.
\end{exercise}

\section{Hints}

\begin{hint}{ex:admissible2}
Only the primes up to the size of the set matter. Look modulo $3$ and $5$.
\end{hint}
\begin{hint}{ex:twin12}
Work modulo $4$ and modulo $3$.
\end{hint}
\begin{hint}{ex:averagegap}
Sum the gaps of the primes in $]\frac x2,x]$.
\end{hint}
\begin{hint}{ex:factorial}
Stirling's formula gives $\log n!\sim n\log n$.
\end{hint}
\begin{hint}{ex:diameter}
The $k$-th prime is about $k\log k$, and so is the $(\pi(k)+k)$-th.
\end{hint}
\begin{hint}{ex:hphiavg}
Write $\frac n{\varphi(n)}=\sum_{d\mid n}\frac{\mu^2(d)}{\varphi(d)}$.
\end{hint}
\begin{hint}{ex:triplesieve}
Now $\rho(p)=3$ for $p\ge5$, $\rho(2)=1$, $\rho(3)=2$; split $d$ into three coprime factors.
\end{hint}
\begin{hint}{ex:jacobsthal}
For $210$, try the integers around $\pm$ a well-chosen multiple of $210$ and use symmetry $n\mapsto-n$.
\end{hint}
\begin{hint}{ex:pigeon}
There are only $123$ even numbers up to $246$.
\end{hint}
\begin{hint}{ex:bruntail}
Abel summation, as in Tome~5, Corollary~8.6.
\end{hint}

\section{Solutions}

\begin{solution}{ex:admissible2}
\idea A set of $k$ elements can only fail at a prime $p\le k$, and fails when it hits every class.
\writeup $\{0,2,4\}$ hits $0,2,1$ modulo $3$: not admissible. $\{0,2,6\}$ is $0,2,0$ modulo $3$ and all even: admissible. $\{0,4,6\}$ is $0,1,0$ modulo
$3$: admissible. $\{0,2,6,8\}$ is $0,2,0,2$ modulo $3$, and only $p\le4$ matter: admissible (the prime quadruplets $5,7,11,13$ and $11,13,17,19$).
$\{0,2,8\}$ is $0,2,2$ modulo $3$: admissible. $\{0,6,12,18,24\}$ is $0,1,2,3,4$ modulo $5$: not admissible.
\end{solution}

\begin{solution}{ex:twin12}
\idea Twins straddle a multiple of $6$, and $p(p+2)=(p+1)^2-1$.
\writeup Both are odd, so $p+1$ is even; among $p,p+1,p+2$ one is divisible by $3$, and it must be $p+1$. So $6\mid p+1$ and $12\mid2(p+1)=2p+2$. And
$p(p+2)+1=(p+1)^2$.
\end{solution}

\begin{solution}{ex:averagegap}
\idea An average is squeezed between the smallest and the largest value.
\writeup Let $q_1<\dots<q_N$ be the primes in $]\frac x2,x]$; $N\sim\frac x{2\log x}$. The $N-1$ gaps sum to $q_N-q_1\le\frac x2$, and each $q_i$ has
$\log q_i\ge\log x-1$. So the smallest gap is at most $\frac x{2(N-1)}=(1+o(1))\log x$, which is $(1+o(1))\log q_i$; this gives the $\liminf$. The gaps
including the one from $q_N$ to the next prime after $x$ sum to at least $q_{N}-q_1\ge\frac x2(1-o(1))$ (by the prime number theorem, there are primes
in $]\frac x2,\frac x2(1+\epsilon)]$ and $]x(1-\epsilon),x]$), so the largest is at least $(1-o(1))\log x$; this gives the $\limsup$.
\end{solution}

\begin{solution}{ex:factorial}
\idea The run has length $n$, but it lives near $n!$, where the logarithm is $n\log n$.
\writeup The prime $p$ before $n!+2$ satisfies $p\le n!+1$, and $p\ge\frac{n!}2$ by Bertrand's postulate, so $\log p=(1+o(1))\log n!=(1+o(1))n\log n$ and
$\log\log p=(1+o(1))\log n$. The gap is at least $n$, and $n=(1+o(1))\frac{\log p}{\log\log p}$. This is much less than the average gap $\log p$.
\end{solution}

\begin{solution}{ex:diameter}
\idea Both ends of the set are primes of index about $k$.
\writeup The set consists of the primes $p_{j}$ with $\pi(k)<j\le\pi(k)+k$, and $\pi(k)=o(k)$. By the prime number theorem $p_j\sim j\log j$, so
$p_{\pi(k)+k}=(1+o(1))k\log k$ and $p_{\pi(k)+1}\le2k$ for large $k$ (Bertrand), and the diameter is $(1+o(1))k\log k$.
\end{solution}

\begin{solution}{ex:hphiavg}
\idea Expand as a divisor sum and swap; the same expansion controls the error terms.
\writeup For any $n$, $\frac n{\varphi(n)}=\prod_{p\mid n}\bigl(1+\frac1{p-1}\bigr)=\sum_{e\mid n}\frac{\mu^2(e)}{\varphi(e)}$. So
$\sum_{n\le x}\frac n{\varphi(n)}=\sum_{e\le x}\frac{\mu^2(e)}{\varphi(e)}\lfloor\frac xe\rfloor=x\sum_{e\le x}\frac{\mu^2(e)}{e\varphi(e)}+O\bigl(\sum_{e\le x}\frac1{\varphi(e)}\bigr)$.
The full series is $S=\sum_e\frac{\mu^2(e)}{e\varphi(e)}=\prod_p\bigl(1+\frac1{p(p-1)}\bigr)$, which converges, and multiplying out gives
$(1-p^{-2})(1-p^{-3})\bigl(1+\frac1{p(p-1)}\bigr)=1-p^{-6}$, so $S=\frac{\zeta(2)\zeta(3)}{\zeta(6)}$.

Both error terms come from the identity $\frac1{\varphi(d)}=\frac1d\sum_{e\mid d}\frac{\mu^2(e)}{\varphi(e)}$. Writing $d=ek$, we get
\[
\sum_{d\le x}\frac1{\varphi(d)}=\sum_{e\le x}\frac{\mu^2(e)}{e\varphi(e)}\sum_{k\le x/e}\frac1k\le S(1+\log x).
\]
Similarly
\[
\sum_{d>x}\frac{\mu^2(d)}{d\varphi(d)}\le\sum_{d>x}\frac1{d^2}\sum_{e\mid d}\frac{\mu^2(e)}{\varphi(e)}=\sum_e\frac{\mu^2(e)}{e^2\varphi(e)}\sum_{k>x/e}\frac1{k^2}\le\frac{2S}x,
\]
using $\sum_{k>y}\frac1{k^2}\le\frac2y$ for every $y>0$. So the main term is $Sx+O(1)$ and the total error is $O(\log x)$.
\end{solution}

\begin{solution}{ex:triplesieve}
\idea Three linear factors give $\rho(p)=3$, and $3^{\omega(d)}$ counts splittings into three coprime factors.
\writeup Sieve $\mathcal A=\{n(n+2)(n+6):n\le x\}$ with $y=z=\frac{\sqrt x}{\log^4x}$. Here $\rho(2)=1$, $\rho(3)=2$ (classes $0$ and $1$), and $\rho(p)=3$ for
$p\ge5$, so $\eta(p)=\frac3{p-3}\ge\frac3{p-1}$ for $p\ge5$. Keeping $d$ prime to $6$, $G(y)\ge\sum\frac{\mu^2(d)3^{\omega(d)}}{\varphi(d)}\ge\sum_{a,b,c\le y^{1/3}}\frac{\mu^2(a)\mu^2(b)\mu^2(c)}{\varphi(a)\varphi(b)\varphi(c)}$
over pairwise coprime $a,b,c$ prime to $6$, which by (i) and (ii) of the proof of Proposition~\ref{prop:pairs}, applied twice, is $\gg\log^3y$. The
remainder is at most $\bigl(\sum_{d\le y}\tau_3(d)\bigr)^2\le y^2(1+\log y)^4$, which is $O(\frac x{\log^4x})$. So the count is $O(\frac x{\log^3x})$.
\end{solution}

\begin{solution}{ex:jacobsthal}
\idea Integers sharing a factor with $m$ are those not coprime to it, and the pattern is periodic and symmetric under $n\mapsto m-n$.
\writeup Modulo $30$, the integers coprime to $30$ are $1$, $7$, $11$, $13$, $17$, $19$, $23$ and $29$; the longest gap between consecutive ones is $6$ (from $1$ to $7$, and from
$23$ to $29$), so the longest run of integers sharing a factor with $30$ has length $5$, for instance $2,3,4,5,6$. For $210$, a check of one period shows
the longest run has length $9$, for instance $2,3,\dots,10$ (the next integer coprime to $210$ after $1$ is $11$); in Jacobsthal's notation $j(210)=10$.

This run contains the primes $2,3,5,7$, so it is not yet a run of composites. Reflecting by $n\mapsto210-n$, which preserves $\gcd(n,210)$, gives the run
$200,201,\dots,208$, and now every term shares a factor with $210$ and exceeds $7$, so every term is composite. That is the argument of
Theorem~\ref{thm:westzynthius} in miniature: $9$ consecutive composites below $420$ (in fact below $210$).

It is not the first such run. The prime gaps below $200$ are at most $8$ (the gap $89\to97$ gives the $7$ composites $90,\dots,96$), and the first gap of
length at least $10$ is $113\to127$. So the first run of $9$ consecutive composites is $114,\dots,122$, inside a run of $13$. Curiously, $114,\dots,126$ is
also the longest run of integers sharing a factor with $2310$ in the first period.
\end{solution}

\begin{solution}{ex:pigeon}
\idea Infinitely many objects in finitely many boxes.
\writeup Every $d_n\ge2$ for $n\ge2$ is even. Infinitely many $d_n$ lie in the $123$ even values $2,4,\dots,246$, so one value occurs infinitely often.
Nobody knows which: proving it for a specific $h$, such as $h=2$, is the whole difficulty.
\end{solution}

\begin{solution}{ex:bruntail}
\idea Partial summation turns a counting bound into a bound for a sum of reciprocals.
\writeup Let $\pi_2(t)\le\frac{Ct}{\log^2t}$ for $t\ge2$. By Abel summation, $\sum_{x<p\le X}\frac1p=\frac{\pi_2(X)}X-\frac{\pi_2(x)}x+\int_x^X\frac{\pi_2(t)}{t^2}\dd t$
over twin primes $p$. Letting $X\to\infty$, the tail is at most $\int_x^\infty\frac C{t\log^2t}\dd t=\frac C{\log x}$.
\end{solution}

\section*{Sources and further reading}

\begin{itemize}
\item A.~de Polignac, \emph{Recherches nouvelles sur les nombres premiers}, Comptes Rendus de l'Acad\'emie des Sciences 29 (1849), 397--401 and
738--739 (the epigraph, paraphrased).
\item G.~H.~Hardy and J.~E.~Littlewood, \emph{Some problems of `Partitio numerorum' III: on the expression of a number as a sum of primes}, Acta
Mathematica 44 (1923), 1--70 (the $k$-tuples conjecture).
\item P.~Erd\H{o}s, \emph{The difference of consecutive primes}, Duke Mathematical Journal 6 (1940), 438--441.
\item E.~Westzynthius, \emph{\"Uber die Verteilung der Zahlen, die zu den $n$ ersten Primzahlen teilerfremd sind}, Commentationes Physico-Mathematicae
Helsingfors 5 (1931), 1--37. R.~A.~Rankin, \emph{The difference between consecutive prime numbers}, Journal of the London Mathematical Society 13 (1938),
242--247.
\item K.~Ford, B.~Green, S.~Konyagin, J.~Maynard and T.~Tao, \emph{Long gaps between primes}, Journal of the American Mathematical Society 31 (2018),
65--105.
\item D.~A.~Goldston, J.~Pintz and C.~Y.~Y\i ld\i r\i m, \emph{Primes in tuples I}, Annals of Mathematics 170 (2009), 819--862.
\item Y.~Zhang, \emph{Bounded gaps between primes}, Annals of Mathematics 179 (2014), 1121--1174. J.~Maynard, \emph{Small gaps between primes}, Annals of
Mathematics 181 (2015), 383--413. D.~H.~J.~Polymath, \emph{Variants of the Selberg sieve, and bounded intervals containing many primes}, Research in the
Mathematical Sciences 1 (2014), article 12 (the bound $246$).
\item A.~Granville, \emph{Primes in intervals of bounded length}, Bulletin of the American Mathematical Society 52 (2015), 171--222; \emph{Harald Cram\'er
and the distribution of prime numbers}, Scandinavian Actuarial Journal (1995), 12--28 (the constant $2e^{-\gamma}$).
\item T.~Oliveira e Silva, S.~Herzog and S.~Pardi, \emph{Empirical verification of the even Goldbach conjecture and computation of prime gaps up to
$4\cdot10^{18}$}, Mathematics of Computation 83 (2014), 2033--2060, and the table of gaps on Oliveira e Silva's web page (the gap $1476$).
\item The counts of twins, cousins and triples up to $10^9$ and the maximal gaps up to $10^9$ were computed for this book with a sieve; the notebook of
this chapter repeats the computation.
\end{itemize}
